% fixture: amsmath/eqref
% tier: T3
% engine: document
% covers: amsmath's \eqref: \textup{\tagform@{\ref{key}}}, the reference in parentheses in an \hbox in \normalfont (with \@@italiccorr), before and after its equation, an undefined key, in italic and bold text, followed by punctuation
\documentclass{article}
\usepackage{amsmath}
\begin{document}
As \eqref{eq:later} shows, and as we prove below, the sum is finite.
\begin{equation}
  a^2 + b^2 = c^2 \label{eq:pyth}
\end{equation}
By \eqref{eq:pyth}, and also \eqref{eq:missing}, we get
\begin{equation}
  e^{i\pi} + 1 = 0. \label{eq:later}
\end{equation}
\textit{In italics: \eqref{eq:pyth} and} \textbf{in bold \eqref{eq:later}}. See \eqref{eq:pyth}.
\end{document}
